%=====================================================================
\chapter{Answers to Checkpoints}\label{app:answers}
%=====================================================================
\normalfigures

Answers to the three-question Checkpoint box in each chapter. Attempt them before
reading. Where an answer requires judgement, the reasoning matters more than the
conclusion.

\newcommand{\ansch}[1]{\subsection*{\color{primary}Chapter #1}}
\newenvironment{ans}{\begin{enumerate}[leftmargin=*, itemsep=4pt, topsep=3pt,
  label=\textbf{\arabic*.}]}{\end{enumerate}}

%---------------------------------------------------------------------
\ansch{1 --- What Artificial Intelligence Is}
\begin{ans}
  \item \textbf{Goal} is passed on the description alone: ``minimise total project
        delay'' is an objective, not a procedure. \textbf{Knowledge} is arguably
        passed too --- the task list with durations and dependencies is an
        explicit, inspectable representation that a delivery manager could correct
        without touching the algorithm. \textbf{Search}, \textbf{adaptation} and
        \textbf{explanation} cannot be judged without seeing the implementation: a
        program with this brief might enumerate alternatives or might compute one
        critical path directly, and the description does not say.
  \item It is succeeding in \emph{acting humanly} --- the bottom-left quadrant, the
        Turing test standard, where indistinguishability from a person is the whole
        criterion. It is failing in \emph{acting rationally}: with no
        representation of the service catalogue it cannot be relied on to give the
        right answer, and with no explanation its answers cannot be checked. It is
        also failing \emph{thinking rationally}, since there is no inference to
        inspect. This is exactly the Eliza pattern of \chref{classic}.
  \item \chref{expert} builds one. Such a system fails the \textbf{adaptation}
        test: its rules are elicited from experts by hand and it does not improve
        with experience. That is precisely the knowledge acquisition bottleneck
        which ended the expert-system boom --- and note that it fails the one test
        this course does not teach, while passing the other four.
\end{ans}

%---------------------------------------------------------------------
\ansch{2 --- Intelligent Agents and Task Environments}
\begin{ans}
  \item A reasonable answer: \emph{Performance measure} --- total time pull
        requests spend waiting, weighted by urgency, plus a penalty for defects
        that reach the main branch. Note that this describes the state of the
        repository. Measures such as ``number of reviews performed'' describe the
        agent and would be gamed by skimming. \emph{Environment} --- the open pull
        requests, their size, age and author, the reviewers and their load, the
        build status. \emph{Actuators} --- assign a pull request to a reviewer,
        defer it, request changes. \emph{Sensors} --- the repository API: diffs,
        timestamps, build results, reviewer availability.
  \item Fully observable (the coins inserted are the whole relevant state),
        deterministic, episodic (each sale is independent), static, discrete, and
        single-agent. All six are the easy case, which is why a \textbf{simple
        reflex agent} is sufficient --- and why a vending machine is not usually
        described as intelligent.
  \item Yes, it was rational. Rationality is maximising \emph{expected} performance
        given the percept sequence available at the time, not achieving the best
        outcome. The agent acted on strong evidence; the outcome was unlucky. This
        distinction is what makes a blameless incident review possible: you judge
        the decision procedure against what was knowable, and you change the
        procedure only if the evidence was misread.
\end{ans}

%---------------------------------------------------------------------
\ansch{3 --- Representations and Search in Python}
\begin{ans}
  \item A list is mutable, and Python makes mutable built-in types unhashable,
        because an object whose hash could change while it sits in a set would
        become unfindable. Without a usable explored set the search cannot ask
        ``have I seen this state?''. Every state is then re-expanded once for each
        distinct path that reaches it, so the search explores the search
        \emph{tree} rather than the state \emph{graph} --- and on a space with a
        cycle that tree is infinite, so the search never terminates.
  \item The heap entries are tuples such as \texttt{(cost, node)}. When two costs
        are equal Python moves on to compare the second elements, and \texttt{Node}
        defines no ordering, so it raises \texttt{TypeError}. It is intermittent
        because it needs a tie, which most inputs do not produce. The fix is one
        element: push \texttt{(cost, next(counter), node)} with a monotonically
        increasing counter, so ties are broken by insertion order and two nodes are
        never compared.
  \item It multiplies the state space by the number of distinct step counts at
        which each situation can be reached. Two searches that arrive at the same
        set of finished tasks after $7$ and after $9$ steps now hold
        \emph{different} states, so the explored set never matches, nothing is
        pruned, and the search degenerates towards enumerating paths. The rule is
        in the chapter: if the goal test cannot see it, it does not belong in the
        state.
\end{ans}

%---------------------------------------------------------------------
\ansch{4 --- Problem Formulation and Uninformed Search}
\begin{ans}
  \item They are repeated states: the same set of finished tasks reached by
        different orderings. The search tree contains one node per \emph{path}, the
        state graph one node per \emph{situation}. An \term{explored set} --- a set
        of states already generated, tested before pushing --- removes them, and on
        a space containing a cycle it is also what makes the search terminate at
        all.
  \item Because when a node is \emph{generated}, a cheaper path to the same state
        may still be sitting unexpanded in the frontier; only when a node is
        \emph{popped} does the frontier ordering guarantee that nothing cheaper
        remains. The smallest graph showing it: start $S$, goal $G$, and one
        intermediate $A$, with $S \to G$ costing $10$, $S \to A$ costing $1$ and
        $A \to G$ costing $2$. Testing on generation returns the direct path at
        cost $10$; testing on popping returns $S,A,G$ at cost $3$.
  \item Because the number of nodes grows geometrically with depth, so the deepest
        level dominates the total: the nodes above it sum to roughly
        $1/(b-1)$ of it. At $b=10$, $d=5$ the repetition costs about 11\% extra
        work. What it buys is memory: $O(bd)$ instead of $O(b^{d})$ --- $50$ nodes
        rather than $100{,}000$ --- while keeping completeness and the
        shallowest-solution guarantee.
\end{ans}

%---------------------------------------------------------------------
\ansch{5 --- Informed Search and Heuristics}
\begin{ans}
  \item Yes, it is admissible: $23 \le 28$, and the lab verifies
        $h(s) \le h^{*}(s)$ on all $55$ reachable states. A heuristic returning
        $30$ at the root would \emph{not} be admissible, because it exceeds the
        true optimal completion cost of $28$. With it, $A^{*}$ loses its optimality
        guarantee: the node on the optimal path can be buried behind a worse goal
        node in the frontier, and the search returns the worse plan with no
        indication that it has done so.
  \item Because greedy best-first search orders its frontier by $h$ alone and never
        consults $g$. Building the harness costs three days --- an increase in $g$
        --- while reducing the mandatory-work estimate by nothing at all, since the
        harness is not mandatory work. So it looks like pure delay. Any decision of
        the form ``spend more now to spend less later'' is structurally invisible
        to greedy search.
  \item $\max(h_{1}, h_{2})$. It is admissible, because at every state it equals
        one of two values that are each at most $h^{*}$; and it dominates both by
        construction. The \emph{average} is also admissible but is dominated by the
        maximum --- it is never larger and usually smaller, so it expands at least
        as many nodes for the same evaluation cost. Take the maximum, never the
        mean.
\end{ans}

%---------------------------------------------------------------------
\ansch{6 --- Local Search and Optimisation}
\begin{ans}
  \item A \term{plateau} is a flat region of the landscape; a \term{shoulder} is a
        plateau with rising ground beyond it. (a)~Forbidding sideways moves means
        halting on every shoulder, so solutions that lie just beyond flat ground
        are never reached. (b)~Allowing them without limit means the search can
        circle a genuine plateau forever. The standard compromise is a cap ---
        allow sideways moves up to a fixed count, then stop.
  \item With $\Delta = -6$: at $T = 8$, $p = e^{-6/8} = e^{-0.75} \approx 0.472$,
        so a move six points worse is accepted almost half the time. At $T = 0.5$,
        $p = e^{-12} \approx 6.1 \times 10^{-6}$, which is effectively never. Early
        in the run the search wanders and can cross a ridge; late in the run it is
        hill climbing and only consolidates.
  \item Crossover assumes the state decomposes into \emph{blocks that are
        independently good}, so that splicing part of one parent onto part of
        another preserves value. The staffing objective largely satisfies this:
        most of a developer's contribution is her own skill score, which depends on
        her assignment alone. What would break it is an objective depending on
        \emph{pairs} of positions --- who is rostered alongside whom, or two
        services that must not share a host --- because then a block's value
        depends on the block it is spliced next to.
\end{ans}

%---------------------------------------------------------------------
% Chapters 7-22: answers are written with each chapter.
%---------------------------------------------------------------------
\ansch{7 --- Constraint Satisfaction}
\begin{ans}
  \item Because nothing in plain backtracking ever shrinks a domain. Every
        unassigned variable still has all four sprints available, so they all tie
        on the minimum-remaining-values count and the heuristic falls back on the
        first variable in the list --- which is exactly what plain backtracking was
        already doing. The change that makes MRV worth having is \emph{propagation}:
        forward checking removes values from neighbours after each assignment, and
        from that moment the domains have different sizes and MRV has something to
        compare. Measured: $3{,}784$ checks with MRV alone, $415$ with MRV and
        forward checking together.
  \item Both are true because AC-3 works only on \emph{binary} constraints. The
        precedence and same-developer constraints are binary, and AC-3 prunes them
        thoroughly --- every one of the twelve domains loses at least one sprint.
        But the constraint that actually binds this problem is capacity: $45$
        points into $48$, an $n$-ary constraint over all twelve variables at once,
        which no arc can express and AC-3 never examines. So the preprocessing
        removes values that the search was not going to waste much time on anyway,
        at a cost of $224$ extra consistency checks and a saving of zero nodes.
  \item Either: instrument the propagation to record which constraint emptied a
        domain most often, and report that one; or re-run the solver with each
        constraint (or constraint class) relaxed in turn and report which
        relaxations make the problem solvable. Here you would expect
        \emph{capacity} to bind, and the lab confirms it --- raising the cap from
        $12$ to $13$ takes the solution count from $14$ into the hundreds, while
        the precedence constraints have plenty of slack.
\end{ans}

\ansch{8 --- Adversarial Search and Game Playing}
\begin{ans}
  \item Because alpha--beta only ever declines to evaluate leaves that
        \emph{provably} cannot change the answer: a cutoff happens exactly when the
        player above already has a better alternative and will therefore never
        choose the branch being examined. The two searches must agree on every
        instance, not just this one --- pruning is exact, and what it approximates
        is nothing at all.
  \item First, the right-hand MIN node has already found a child worth $2$, so it
        can hold us to at most $2$ whatever its remaining child contains. Second,
        the root has already secured $5$ from the left subtree. Since $2 < 5$ we
        will never play into this branch, so the value of its unexamined child
        cannot affect the root --- and alpha--beta never learns it.
  \item The probes are themselves evaluations, and on a shallow tree they cost more
        than the cutoffs they buy: the lab measures $506$ leaf evaluations with a
        depth-1 probe ordering against $75$ with natural ordering. For ordering to
        pay, the tree must be deep enough that a better ordering saves an
        exponential amount, and the probe results must be cached --- a
        transposition table --- so that the work done to order one search is reused
        by the next. That is the situation in a chess engine at twenty ply, and not
        the situation at five.
\end{ans}

\ansch{9 --- Planning}
\begin{ans}
  \item \texttt{test A} has \texttt{staging free} in its \emph{delete} list, so
        testing the first component takes the staging environment away. But
        \texttt{test B} has \texttt{staging free} in its \emph{preconditions}, so
        the second component cannot be tested until something restores it, and
        \texttt{release staging} is the only operator that does. It is there purely
        to undo interference caused by an earlier step. Removing
        \texttt{staging free} from \texttt{test}'s delete list makes it disappear:
        actions can then no longer interfere, and the problem collapses into the
        topological sort of \chref{uninformed}.
  \item Because $A^{*}$ saves work only by refusing to expand nodes with
        $f(n) \ge C^{*}$, and in this domain almost every reachable state lies on
        \emph{some} optimal plan --- nearly every action is necessary and their
        order is largely free --- so $f$ is close to $C^{*}$ nearly everywhere and
        there is almost nothing to rule out. The same heuristic would pay well on a
        planning problem with many \emph{dead ends or wasted actions}: irreversible
        mistakes, resources that can be spent on the wrong thing, or a much larger
        operator set of which only a few are relevant. Heuristics pay in proportion
        to how much of the state space is wasted effort.
  \item Completeness. Breadth-first progression search will find a plan whenever one
        exists; means--ends analysis takes the goals in some order and can destroy,
        while achieving an early goal, a condition a later goal needs. The two-operator
        example of \dsref{clobber} is solvable, and means--ends analysis reports
        failure on it when the goals are taken in one of the two possible orders.
        (Breadth-first search also guarantees a \emph{shortest} plan; means--ends
        analysis found optimal plans here but is not obliged to.)
\end{ans}

\ansch{10 --- Knowledge Representation}
\begin{ans}
  \item The walk starts at \texttt{Release 4.7}, finds no local
        \texttt{regression-pass}, moves to \texttt{Hotfix}, finds \texttt{smoke},
        and \emph{stops}. \texttt{Release} is never consulted, so its \texttt{full}
        never competes. ``Most specific wins'' does not have to be stated as a rule
        because it is a consequence of the search order: the walk goes from
        specific to general and halts at the first hit.
  \item \term{Non-monotonicity} --- adding knowledge removed a conclusion.
        Classical logic does not have it: entailment is defined as truth in every
        model of the knowledge base, and adding a sentence can only rule models
        out, so the set of entailed conclusions can only grow. That is exactly the
        trade \chref{propositional} makes --- it gains soundness and loses the
        ability to say ``normally''.
  \item A semantic network with defaults can express \emph{normally}: ``releases go
        out on Thursdays'' with hotfixes excepted, without the exception making the
        knowledge base inconsistent. Written as a universally quantified sentence
        that claim is simply false, and adding one hotfix makes the knowledge base
        contradictory. What logic can do and the network cannot is \emph{derive
        guaranteed consequences}: a sound procedure over logic gives conclusions
        that hold in every model, whereas an inheritance walk gives whatever the
        nearest ancestor happened to say, with no guarantee it is consistent with
        anything else in the network.
\end{ans}

\ansch{11 --- Propositional Logic and Inference}
\begin{ans}
  \item $\mathit{KB} \models \alpha$ is about \emph{meaning}: $\alpha$ is true in
        every model in which the knowledge base is true, and no algorithm is
        mentioned. $\mathit{KB} \vdash \alpha$ is about a \emph{procedure}: that
        procedure can produce $\alpha$. A procedure that derives every sentence is
        \emph{complete} --- it certainly derives everything entailed --- but wildly
        \emph{unsound}, since it derives a great many things that are not entailed.
        Soundness is the property it violates, and it is the one that matters:
        completeness alone is worthless.
  \item F1--F3 force $\mathit{Reviewed}$, $\mathit{Tested}$ and
        $\mathit{SignedOff}$; R1 then forces $\mathit{Staged}$; R2 then forces
        $\mathit{Deploy}$. R3 says $\mathit{Blocker} \Rightarrow \neg
        \mathit{Deploy}$, and since $\mathit{Deploy}$ is true, $\mathit{Blocker}$
        must be false or R3 is violated. Exactly one model survives, so both
        $\mathit{Deploy}$ and $\neg\mathit{Blocker}$ are entailed.
        It is \emph{not} a bug in the inference procedure --- the procedure drew a
        correct consequence of what it was given. It is a bug in the \emph{axioms}:
        the team meant ``deployment is permitted provided there is no blocker'',
        and R2 as written does not say that. Adding $\neg\mathit{Blocker}$ to R2's
        antecedent makes the entailment disappear, which the lab demonstrates.
  \item First, resolution produces an \emph{auditable proof} --- a numbered list of
        clauses each derived from two earlier ones, which a person can check line by
        line and an auditor will accept. A model count proves nothing a human can
        read. Second, and decisively, resolution \emph{lifts to first-order logic}
        (\chref{fol}), where the domain may be infinite, there are infinitely many
        models, and enumerating them is not slow but impossible. DPLL's speed is
        confined to the propositional case.
\end{ans}

\ansch{12 --- First-Order Logic and Resolution}
\begin{ans}
  \item $\forall x \; \mathit{service}(x) \Rightarrow \exists y \;
        \mathit{owns}(y, x)$. Eliminating $\Rightarrow$ gives
        $\forall x \; \neg\mathit{service}(x) \vee \exists y \;
        \mathit{owns}(y,x)$; skolemising $y$ and dropping the universal gives
        $\neg\mathit{service}(x) \vee \mathit{owns}(\mathit{owner}(x), x)$.
        The existential variable lies \emph{inside} the universal quantifier over
        $x$, so the owner depends on which service we are talking about. Skolemising
        to a constant \texttt{Sk} would give
        $\neg\mathit{service}(x) \vee \mathit{owns}(\texttt{Sk}, x)$ --- ``one
        particular person owns every service'' --- which is a different and much
        stronger claim.
  \item $n^{2}$ sentences: one for each ordered pair of services, since the rule
        relates a depender to a dependee. Adding one service means writing $2n$ more
        by hand and getting none of them wrong. In first-order logic the rule is a
        single sentence and adding a service means adding a \emph{fact}; the rule is
        untouched. That is the difference between a knowledge base a team can
        maintain and one it cannot.
  \item Because first-order entailment is only \term{semi-decidable}. If the
        conclusion follows, resolution will eventually derive the empty clause; if
        it does not, the procedure may search forever, and it cannot in general
        report that no proof exists. A countermodel would require exhibiting an
        interpretation in which the knowledge base holds and the query fails, and
        first-order interpretations can have infinite domains --- there are
        infinitely many to consider and no way to enumerate them. Propositional
        logic has finitely many models, which is exactly why
        \chref{propositional} could hand one back.
\end{ans}

\ansch{13 --- Rule-Based Expert Systems}
\begin{ans}
  \item \texttt{wake\_budget}: \texttt{available} for T-4471 and \texttt{spent} for
        T-4473. It is \emph{not} a fact about the ticket --- it is a fact about the
        on-call engineer and what has already been asked of her tonight. Both
        tickets reach \texttt{escalate = yes}; they differ only at R14 against R15,
        which consult that fact. This is precisely the internal state
        \chref{agents} showed a model-based reflex agent must maintain because no
        percept reports it, and here it costs one entry in working memory and one
        pair of rules.
  \item The nine is right for a monitoring dashboard or an alerting pipeline, where
        a newly arrived fact should trigger \emph{everything} that follows from it
        --- notify finance, attach the payments runbook, set the page severity ---
        without anybody having to ask for each one. The six is right when an
        engineer or an auditor has a single question (``why was this escalated?''),
        because forward chaining would derive a great deal that is irrelevant to it
        and would still not produce a proof tree.
  \item They get the backward-chaining explanation: \texttt{escalate = no} by R11
        because \texttt{needs\_owner = no} by R7, because the priority is P2 by R3
        (major severity, standard tier) and there was a workaround. Every step names
        the rule and the fact it used, so the auditor can check each line and, if
        they disagree, point at the specific rule that is wrong. A probability
        satisfies none of this: it cannot be checked, it identifies nothing to
        correct, and it does not say which of the ticket's properties mattered. See
        also \chref{responsible}, where this is argued as accountability rather than
        convenience.
\end{ans}

\ansch{14 --- Uncertainty and Decision Making}
\begin{ans}
  \item Of a thousand releases, a hundred are defective and the test catches ninety
        of them --- but nine hundred are sound and the test cries wolf on a fifth of
        those, which is a hundred and eighty false alarms. So of the two hundred and
        seventy alarms, only ninety are real. The test is good; there are simply far
        more sound releases for it to be wrong about, and that is the base rate
        doing the work.
  \item The lesson is that \emph{a change of belief is not a change of decision}.
        The recommendation depends only on which side of the flip point $p^{*}$ the
        probability falls, and $1/3$ is below $10/29 \approx 0.345$. What an
        engineer should report is the \emph{action} and the \emph{margin}: ``still
        ship, but only by $1.67$ utility units, which is close enough that better
        evidence is now worth having'' --- which is both the decision and the case
        for the next step.
  \item Because the value of information is proportional to how likely it is to
        \emph{change the action}, not to how interesting it is. Before the test the
        decision was comfortable --- ship, by $35$ units --- so certainty could
        rarely have flipped it. After the alarm the two options differ by under two
        units, so the same certainty now flips the decision often. Information is
        worth most exactly where the decision is hardest, and the function peaks at
        the flip point itself.
\end{ans}

\ansch{15 --- Bayesian Networks}
\begin{ans}
  \item It asserts that \emph{each variable is conditionally independent of its
        non-descendants given its parents}. That licenses the chain rule, so the
        joint factorises into one small table per node. An assertion in
        \dsref{bayesnet} that could easily be false: there is no edge between
        \emph{Rushed} and \emph{Junior}, which claims they are independent. In many
        organisations rushed work is precisely what gets handed to the least
        experienced developer, and if that is true here the network is wrong and
        will give confidently wrong answers.
  \item Both are consequences of the arithmetic of a common effect.
        \emph{Defect} needs explaining, and a rushed specification is one of the
        available explanations, so its probability rises. Learning that a junior
        developer was involved supplies a \emph{different} sufficient explanation,
        so less explanatory work is left over and the probability of the first cause
        falls. A forward-chaining rule base cannot reproduce this because it would
        need a rule whose conclusion becomes \emph{weaker} when a condition is
        added --- and \chref{expert}'s specificity strategy would do the opposite,
        preferring the rule with more conditions and thus strengthening it.
  \item It becomes useless when the evidence is rare: it keeps only the fraction of
        samples that happen to match, so evidence of probability $0.01$ throws away
        $99\%$ of the work, and diagnostic evidence is usually rare because that is
        what makes it worth reasoning about. Likelihood weighting does not sample
        the evidence variables at all --- it fixes them to their observed values and
        weights each sample by the probability of the evidence given what was
        sampled --- so every sample contributes.
\end{ans}

\ansch{16 --- Fuzzy Logic}
\begin{ans}
  \item Because a membership is a fact about the \emph{word}, not about our
        knowledge of the project. We know exactly how big the project is: seventy
        person-days. The number $0.4$ says that seventy days is ``high effort'' to
        degree $0.4$, and no further measurement of the project can change that.
        What \emph{would} change it is redefining the membership function --- that
        is, agreeing with the team on what ``high'' should mean, which is a
        conversation about vocabulary rather than about evidence.
  \item The membership functions form a \term{fuzzy partition}: overlapping
        triangles whose memberships sum to one everywhere, so at most two words are
        active for each input and therefore at most $2 \times 2 = 4$ rules fire.
        If the functions do not overlap, only one word is active per input, only one
        rule fires, and there is no interpolation at all --- the controller
        degenerates into a lookup table with the smoothness that motivated the whole
        approach thrown away.
  \item Centroid defuzzification returns the centre of area of the aggregated shape.
        Even when only the lowest output set is active, that set is a triangle
        centred at zero with area extending to its right, so its centre of area is
        strictly positive; the same argument at the top end keeps it below the
        maximum. The consequence for reporting is that the number must never be
        presented as a percentage or as an absolute risk. It is an ordinal score,
        useful for comparing projects with each other, and a board told ``$83\%$
        risk'' has been misled by an artefact of the defuzzifier.
\end{ans}

\ansch{17 --- Learning from Examples}
\begin{ans}
  \item Because the rule learns \emph{only from mistakes}: when the prediction is
        correct, $y - \hat{y} = 0$ and every weight is unchanged. Ten of the twelve
        examples were already classified correctly when they were seen, so they
        contributed nothing. The second update changed the meaning of the rule: it
        drove the weight on clarity negative, so that higher clarity now argues
        \emph{against} lateness. Nobody told the perceptron that; it followed from a
        single misclassified example in the second pass.
  \item Either the data are separable and it simply needs more epochs, or no line in
        the hypothesis class fits the data at all --- and the two traces look
        identical, because in both cases the weights are still moving when the
        budget expires. The cheapest experiment that distinguishes them is a brute
        force sweep over the hypothesis class: on a small problem, enumerate a grid
        of lines and record the best accuracy any of them achieves. On the
        exclusive-or arrangement that answer is $75\%$, which settles the question;
        no learning rate or initialisation can beat a bound over the whole class.
  \item A single perceptron cannot represent exclusive-or, because it computes a
        linear threshold and no straight line separates those four points. The 1969
        conclusion was that neural network research was a dead end, and funding
        largely stopped for fifteen years. It was wrong because the limitation
        applies to a \emph{single layer}: a network with one hidden layer represents
        exclusive-or easily, and what was actually missing was a way to \emph{train}
        such a network, supplied by backpropagation in the mid-1980s. A correct
        result about the simplest case was read as a verdict on the approach.
\end{ans}

\ansch{18 --- Natural Language Processing}
\begin{ans}
  \item Reading~1 attaches \texttt{after the timeout} to the verb phrase:
        \emph{restart the server, and do it after the timeout occurs}. Reading~2
        attaches it to the noun phrase: \emph{restart the server that handles
        timeouts}. They are different instructions. A routing system should not
        choose --- it should recognise that the subject line supports both and ask
        the reporter, which is a capability the grammar provides for free and which
        a system returning one confident parse throws away.
  \item That fluency is not evidence of understanding. Eliza had no representation
        of the conversation, no model of the person and no goal, and it still
        produced responses people found so apt that they disclosed things to it they
        would not tell a colleague --- including people who had read its source. So
        your impression of a transcript tells you about your own response, not about
        the system, and the instrument to use instead is
        \chref{introduction}'s checklist: is there a goal, a search, an inspectable
        representation, adaptation, an explanation?
  \item Sparsity is the problem that most possible n-grams never appear in any
        finite corpus, so unseen ones receive probability zero. It is an estimation
        problem, and smoothing solves it by reserving probability mass for the
        unseen. The fixed window is a \emph{representational} limit: in
        \emph{the server that the engineer who wrote the runbook maintained was
        restarted}, the subject of \emph{was restarted} is nine words back. No
        $n$-gram with a practical $n$ can span that, and raising $n$ makes sparsity
        exponentially worse --- so the two limitations pull against each other and
        no corpus size resolves the second.
\end{ans}

\ansch{19 --- Classic AI Systems and How to Analyse Them}
\begin{ans}
  \item GPS is means--ends analysis (\chref{planning}); Eliza is pattern
        substitution (\chref{nlp}); Student is pattern matching into equations,
        which is \chref{fol}'s unification restricted to strings; Macsyma is term
        rewriting, which is \chref{expert}'s forward chaining applied to expression
        trees rather than to facts.
  \item Because GPS's failure is \emph{legible}: it stops and says it has no plan,
        so a person knows to look. Eliza always has a reply, so nothing distinguishes
        a case it handled well from one it did not handle at all, and the user has
        no signal to act on. The consequence for evaluating a modern language system
        is direct: it too always produces output, so the fluency of a response
        carries no information about whether it is grounded, and evaluation must
        come from an external check --- a test suite, a type checker, a retrieved
        source --- rather than from reading the answer.
  \item A person, writing it down by hand. It predicts that the system will work
        well in the region somebody took the trouble to encode and will need a
        person for every extension --- so as the business changes, the system's
        coverage falls behind at whatever rate knowledge engineers can be hired and
        briefed. That is the knowledge acquisition bottleneck, and it is the reason
        all four of these systems have the biographies they do.
\end{ans}

\ansch{20 --- Recent Trends and Applications}
\begin{ans}
  \item Because \chref{expert}'s chain is a \emph{record} of the inference that
        actually occurred: each line names the rule that fired and the facts it
        matched, and the conclusion was produced by exactly those steps. A language
        model's explanation is generated by the same next-token process as the
        answer itself --- it is text that resembles a derivation, produced
        independently of whatever computation led to the answer. In
        \chref{propositional}'s terms, it may display something shaped like
        $\vdash$ while carrying no guarantee of any relation to $\models$; a
        plausible-looking derivation is not a derivation.
  \item The other five are the loop itself, the tools, the step budget, the
        termination test and the state the agent keeps between steps. Removing the
        \emph{policy's handling of a failed tool call} causes the commonest deployed
        failure: a tool returns nothing, no rule notices, and the \texttt{None}
        travels until it reaches arithmetic and crashes --- or, with a policy that
        cannot terminate, the loop spins until the step budget stops it. Both
        failures live in the engineering around the model, not in the model.
  \item First, \emph{what did the demonstration exclude?} --- every impressive
        demonstration is a microworld until shown otherwise, and the answer tells you
        what the vendor chose not to show. Second, \emph{how would anyone find out
        when it is wrong?} --- the failure-mode column. If the answer is that a
        human notices, ask how they notice when the output is fluent; if it is a
        test or a type checker, the system is safe in proportion to that checker's
        coverage.
\end{ans}

\ansch{21 --- Responsible AI and Its Limits}
\begin{ans}
  \item Through a \term{proxy}: \texttt{workaround} is recorded for about $80\%$ of
        one team's tickets and $15\%$ of the other's, because the first team has
        mature runbooks to cite. R3 makes those tickets P2 and R7 then decides a P2
        with a workaround needs no owner, so that team's tickets are systematically
        not escalated. Deleting an attribute would not have prevented it, and could
        not have: the team name was never an input. The correlation lives in the
        world rather than in a column, which is why \emph{fairness through
        unawareness} is the one approach that reliably fails.
  \item It established \emph{what} produced the gap --- a difference in
        documentation practice being read by the rules as a difference in
        operational need --- and that no rule is individually wrong. It explicitly
        did not establish whether the gap is acceptable. If the first team really
        does need fewer escalations because its runbooks work, the gap is the system
        functioning correctly; if the second team is being under-served because
        nobody has written its runbooks, it is not. That is a management decision,
        and the audit's contribution is to put it in front of somebody who can make
        it, with evidence attached, instead of leaving it to be made silently two
        hundred times a week.
  \item The alternative explanation is \term{automation bias}: the reviewer has
        stopped reviewing and is deferring to the system, which produces a near-zero
        override rate whether the system is right or not. The measurement that
        distinguishes them is to sample decisions for genuine independent review ---
        have a second person decide without seeing the recommendation --- and
        compare. If the independent decisions disagree materially more often than
        the override rate suggests, the low rate was deference rather than
        agreement.
\end{ans}

\ansch{22 --- The Artificial Intelligence Project}
\begin{ans}
  \item It fails the second test: the rules do not exist and would have to be learned
        from historical tickets, which makes it a \emph{Machine Learning} project
        rather than one this course has the machinery for. The advice is to keep the
        domain and change the question to one whose rules can be written down --- for
        example \emph{decide which open tickets to escalate tonight given a stated
        escalation policy} (\chref{expert}), or \emph{schedule this sprint's stories
        without violating capacity or precedence} (\chref{csp}). Both are one
        sentence, both have rules an operations lead can state in an afternoon, and
        both have an embarrassing baseline.
  \item Because a measure chosen after the results are in will be chosen --- honestly
        and without anyone noticing --- to be one the system does well on. The
        failure mode is a report that demonstrates rather than evaluates: it shows
        the system succeeding at whatever it happens to be good at, and provides no
        evidence about the question that was originally asked. Writing the measure
        and the failure condition down in advance is the only practical defence, and
        changing them later is acceptable provided the change and its reason are
        reported.
  \item The baseline it was compared against, the measure it was compared on with
        evidence that the measure was fixed in advance, and the size of the margin
        relative to the variability --- a small margin on one instance with a seeded
        run is not a result. Ideally also the ablation showing which component of
        the system produced the margin, since without it the claim is about the
        whole program rather than about the method being taught.
\end{ans}
